\documentclass[11pt,a4paper]{article}

\usepackage[T1]{fontenc}
\usepackage[utf8]{inputenc}
\usepackage{lmodern}
\usepackage[french]{babel}
\usepackage[a4paper,margin=2.3cm]{geometry}
\usepackage{microtype}
\usepackage{xcolor}
\usepackage{enumitem}
\usepackage{titlesec}
\usepackage{hyperref}
\usepackage{tabularx}
\usepackage{array}
\usepackage{booktabs}
\usepackage[most]{tcolorbox}

% ---- palette (identique à l'étude) ----
\definecolor{ink}{HTML}{1c2330}
\definecolor{accent}{HTML}{2f6e4c}
\definecolor{accent2}{HTML}{8a5a1c}
\definecolor{soft}{HTML}{eef3ef}
\definecolor{softline}{HTML}{cfe0d4}
\definecolor{warn}{HTML}{7a2418}
\definecolor{warnsoft}{HTML}{fbeeea}
\definecolor{warnline}{HTML}{e7c3ba}
\definecolor{open}{HTML}{1c4e7a}
\definecolor{opensoft}{HTML}{eaf1f8}
\definecolor{openline}{HTML}{bcd3e8}

\hypersetup{colorlinks=true, linkcolor=accent, urlcolor=accent2,
  pdftitle={Direction de design --- synthèse d'interview}}

\titleformat{\section}{\Large\bfseries\color{ink}}{\thesection.}{0.5em}{}
\titleformat{\subsection}{\large\bfseries\color{accent}}{\thesubsection}{0.6em}{}
\titlespacing{\section}{0pt}{1.4em}{0.5em}
\setlist[itemize]{leftmargin=1.2em,itemsep=2pt,topsep=2pt}

% ---- callouts ----
\newtcolorbox{principe}{enhanced,breakable,colback=soft,colframe=accent,
  boxrule=0.4pt,left=8pt,right=8pt,top=5pt,bottom=5pt,arc=2pt,
  fonttitle=\bfseries,coltitle=accent,title=Le fil rouge}
\newtcolorbox{antipat}{enhanced,breakable,colback=warnsoft,colframe=warn,
  boxrule=0.4pt,left=8pt,right=8pt,top=5pt,bottom=5pt,arc=2pt,
  fonttitle=\bfseries,coltitle=warn,title=Résolu}
\newtcolorbox{ouvert}{enhanced,breakable,colback=opensoft,colframe=open,
  boxrule=0.4pt,left=8pt,right=8pt,top=5pt,bottom=5pt,arc=2pt,
  fonttitle=\bfseries,coltitle=open,title=Ouvert --- à prototyper}
\newtcolorbox{visionbox}{enhanced,breakable,colback=soft,colframe=accent,
  boxrule=0.8pt,left=10pt,right=10pt,top=7pt,bottom=7pt,arc=3pt}

\newcommand{\jeu}[1]{\textbf{#1}}
\newcommand{\refc}[1]{\hfill{\footnotesize\textcolor{accent2}{#1}}}

\title{\vspace{-1.5em}\textbf{\color{ink}\Huge Direction de design}\\[0.3em]
\large\color{accent}Synthèse d'interview --- comment diriger l'ensemble\\[0.2em]
\normalsize\color{ink}Consolidation des choix de design sur la base de l'étude de progression}
\author{}
\date{}

\begin{document}
\maketitle
\vspace{-3em}

\begin{visionbox}
\textbf{\color{ink}Vision en une ligne.} Un tactique-RPG «~\textbf{XCOM tempéré}~» où \textbf{monter une unité, c'est verrouiller une voie de perks}~; où l'on \textbf{arbitre entre opportunités concurrentes} sur une carte~; et où l'on \textbf{ne perd jamais ses héros --- mais on les épuise}.
\end{visionbox}

\tableofcontents

% ======================================================================
\section{La colonne vertébrale}
\emph{Choix confirmés à l'interview, cohérents entre eux.}

\renewcommand{\arraystretch}{1.3}
\begin{tabularx}{\linewidth}{@{}>{\bfseries}p{3.1cm} X r@{}}
\toprule
\textmd{\textbf{Couche}} & \textbf{Décision} & \textbf{Réf.} \\
\midrule
Moteur & La \textbf{décision / bifurcation} avant l'accumulation & B / C \\
Structure & \textbf{Geoscape concurrent} --- missions simultanées, coût d'opportunité & A3 \\
Pression & \textbf{Horloge douce / locale}, pas d'horloge globale écrasante & A1 \\
Forme & \textbf{Hybride} --- ossature narrative fait-main + variance & A6 \\
Contenu & \textbf{Procédural borné + golden path} & A6 \\
Puissance & \textbf{Niveaux discrets + perks}, lisibles et plafonnés & B3 \\
Profondeur build & \textbf{Perks exclusifs par unité} (le build = portes fermées) & B1 \\
Réversibilité & \textbf{Définitive} --- pas de respec, pas d'undo & --- \\
Économie & \textbf{Monnaies séparées, quasi pas de loot} (effacées) & B8 \\
Mort & \textbf{Asymétrique} --- troupes mortes, héros K.O. & C1 \\
Usure & \textbf{Stress + Fatigue} $\rightarrow$ rotation forcée du roster & C3 \\
Échec & \textbf{Produit du contenu}, jamais de reload & C4 \\
Macro & \textbf{Choix exclusifs réels} --- contenu manquable & C5 \\
Récit & \textbf{Personnel léger} (faits d'armes, historique) & C6 \\
Identité & \textbf{Forte / customisation} --- soldats nommés & B9 \\
\bottomrule
\end{tabularx}
\renewcommand{\arraystretch}{1}

% ======================================================================
\section{Le fil rouge --- ce qui relie tout}

La décision a du \textbf{poids} parce qu'elle est \textbf{irréversible} : perks définitifs, choix macro exclusifs, échec qui fait avancer. Aucune de ces portes ne se rouvre.

La \textbf{rareté} --- moteur classique de la décision --- n'est \textbf{pas} placée dans l'argent ou le loot (effacés), mais dans le \textbf{temps de tes unités}.

\begin{principe}
\textbf{Stress + Fatigue + mort asymétrique} rendent tes meilleures unités régulièrement \textbf{indisponibles}. Tu es donc forcé de \textbf{faire tourner un roster}, donc d'\textbf{investir dans plusieurs builds divergents} (puisqu'ils sont non-respecables).

\smallskip
\textbf{Une seule pression} --- «~mes bonnes unités sont indisponibles~» --- alimente à la fois l'\textbf{enjeu} (jouer sous-optimal) \emph{et} la \textbf{décision} (dans qui j'investis, qui je repose, qui je sacrifie). Macro et micro tirent sur la même corde.
\end{principe}

% ======================================================================
\section{Détail par couche}

\subsection{A --- Macro}
\begin{itemize}
\item \textbf{A3 Geoscape concurrent.} Plusieurs opportunités ouvertes en même temps ; on ne peut pas tout faire $\rightarrow$ coût d'opportunité permanent.
\item \textbf{A1 Horloge douce/locale.} Pression portée par chaque mission/zone (objectifs en N tours, menaces locales), pas par une horloge globale anxiogène.
\item \textbf{A6 Hybride + procédural borné.} Campagne à ossature écrite (golden path) ; cartes/rencontres générées par gabarits bornés autour. L'éditeur de maillage sert à faire-main les moments clés.
\item \textbf{C5 Choix exclusifs réels.} Sauver une zone peut en condamner une autre ; des branches verrouillent du contenu. Assumé : on ne verra pas tout en une partie.
\end{itemize}

\subsection{B --- Micro}
\begin{itemize}
\item \textbf{B3 Niveaux discrets + perks.} Progression lisible, plafonnée ; la puissance vient des \textbf{choix}, pas d'un robinet de stats.
\item \textbf{B1 Perks exclusifs par unité.} À chaque grade, capacité «~A ou B~» irréversible. L'arbre parcouru \textbf{est} l'identité de l'unité.
\item \textbf{Monnaies séparées + quasi pas de loot.} Économie et stuff ne sont \textbf{pas} des leviers de progression. Le poids décisionnel est concentré sur les perks.
\item \textbf{Réversibilité définitive.} Pas de respec, pas d'undo : ce qui donne tout leur poids aux choix de perks. \textbf{Pilier le plus solide du profil.}
\end{itemize}

\subsection{C --- Tension / enjeux}
\begin{itemize}
\item \textbf{C1 Mort asymétrique.} Héros nommés K.O. (récupérables), troupes génériques mortes. Protège l'investissement en builds, garde un enjeu réel.
\item \textbf{C3 Stress + Fatigue.} Deux jauges hors-combat à récupération lente $\rightarrow$ rotation du roster. \textbf{Pas} d'attrition financière (économie effacée).
\item \textbf{C4 L'échec fait avancer.} Un revers monte la menace / ouvre une branche, la partie continue. Tue le save-scum, alimente le récit.
\item \textbf{C6 Récit personnel léger.} Les unités accumulent un historique (faits d'armes, surnoms) ; \textbf{pas} de simulation sociale lourde.
\item \textbf{B9 Identité forte.} Unités nommées et customisables $\rightarrow$ on s'attache, donc le K.O./la rotation portent émotionnellement.
\item \textbf{Pas de séquelles permanentes (C2 écarté).} Le corps des persos reste simple ; l'usure passe par stress/fatigue.
\end{itemize}

% ======================================================================
\section{Frictions avec l'étude --- arbitrage}

\subsection{Friction 1 --- RNG vs moteur «~décision~»}
Hésitation entre \textbf{RNG tempéré} (touches fiables, variance sur dégâts/init) et \textbf{RNG plein} (ratages à \%).
\begin{itemize}
\item Un point déjà fixé tranche en partie : \textbf{pas de save-scum + l'échec fait avancer}. Le gros défaut du RNG plein (rechargement frustré) \textbf{n'existe pas ici} --- un ratage devient une \emph{war story} qui alimente le geoscape. Le RNG plein est donc \textbf{plus défendable} chez toi qu'ailleurs.
\item À l'inverse, le RNG brouille la lisibilité décisionnelle («~bien joué, perdu au dé~»).
\end{itemize}
\begin{ouvert}
Socle = \textbf{RNG tempéré}, avec un \textbf{curseur de variance} assumé. À régler manette en main : c'est \emph{le} point qui se tranche en prototype, pas sur le papier.
\end{ouvert}

\subsection{Friction 2 --- «~mix nombre + level-scaling~» vs anti-pattern}
Choix retenu : \textbf{mix franc (nombre + stats), mais scaling-stats borné par zone} (paliers, pas une fonction continue qui colle au joueur).
\begin{antipat}
\textbf{Pourquoi le borner.} Le treadmill vient du scaling \textbf{vertical en lockstep} : si les stats ennemies montent au même rythme que les tiennes, l'écart relatif reste constant et tu ne te sens jamais plus fort (\emph{l'inflation verticale dévalorise l'effort passé}).

\smallskip
\textbf{Pourquoi ce n'est pas un problème ici.} Ta progression est surtout \textbf{horizontale} (perks = nouveaux \emph{verbes}, pas de plus gros nombres). Un nouveau verbe ne se fait pas avaler par +X\% de PV ennemi. Le scaling-stats borné dose la difficulté \textbf{sans} manger tes gains. Seule la part numérique doit rester bornée.
\end{antipat}

% ======================================================================
\section{Ce que la direction EXCLUT (anti-scope)}
Aussi important que les choix positifs --- ce qu'on s'interdit pour rester cohérent :
\begin{itemize}
\item Pas de loot-grind / robinet à objets (le stuff n'est pas un levier de puissance).
\item Pas de respec / undo / save-scum (les choix sont définitifs).
\item Pas de séquelles/mutilations permanentes.
\item Pas de simulation sociale lourde (récit émergent reste \emph{léger}).
\item Pas d'horloge globale écrasante (pression locale seulement).
\item Pas de level-scaling continu non-borné.
\item Pas de montée verticale comme moteur principal (les stats sont du décor).
\end{itemize}

% ======================================================================
\section{Prochaines étapes --- traduire en pratique}
\begin{enumerate}
\item \textbf{Combat «~Cœur maillage~».} Spécifier la boucle de combat sur le maillage : verbes de base, système de perks (arbre A/B par grade), tour-par-tour, dosage RNG (proto).
\item \textbf{Jauges Stress + Fatigue.} Modéliser les deux jauges, leur récupération à slots limités, l'effet sur la disponibilité du roster.
\item \textbf{Geoscape minimal.} Nœuds concurrents + horloge locale par mission ; brancher l'éditeur de campagne existant.
\item \textbf{Boucle «~échec fait avancer~».} Définir 2--3 conséquences de revers qui produisent du contenu plutôt qu'un game-over.
\end{enumerate}

\begin{center}\small\itshape
Friction 1 (RNG) à trancher au stade prototype combat. Tout le reste est figé.
\end{center}

\end{document}
